\documentclass[11pt]{article}
\usepackage[margin=0.95in]{geometry}
\usepackage{amsmath,amssymb,amsthm,mathtools}
\usepackage{graphicx}
\usepackage{microtype}
\usepackage{enumitem}
\usepackage{xcolor}
\usepackage[colorlinks=true,linkcolor=blue!55!black,citecolor=blue!55!black,urlcolor=blue!55!black]{hyperref}

\newtheorem{theorem}{Theorem}
\newtheorem{lemma}{Lemma}
\newtheorem{remark}{Remark}
\newcommand{\R}{\mathbb R}
\newcommand{\perpB}{\perp_{\!B}}
\newcommand{\Ext}{\operatorname{Ext}}

\title{Three Extreme Points Do Not Force an Isometry\\
\large A full counterexample to the open question in arXiv:2511.20345}
\author{Discovery run \texttt{fa\_banach\_001} -- \texttt{agent\_lane\_18}}
\date{August 11, 2026}

\begin{document}
\maketitle

\begin{center}
\fbox{\parbox{0.92\textwidth}{\textbf{Status: candidate new full counterexample.}
An explicit rational three-dimensional polyhedral norm and an invertible
rational operator give a negative answer. Every orthogonality-preservation
claim reduces to entrywise nonnegativity of one of two displayed $3\times3$
matrices. Independent expert review is nevertheless recommended.}}
\end{center}

\section{Source question}

Manna, Mandal, Paul, and Sain ask on PDF page 16 of \cite{MMPS} whether, in
an $n$-dimensional polyhedral Banach space, preservation of Birkhoff--James
orthogonality at any $n$ linearly independent extreme points forces a bounded
linear operator to be a scalar multiple of an isometry.

\begin{figure}[ht]
  \centering
  \includegraphics[width=0.98\textwidth]{figures/open_problem_crop.png}
  \caption{The exact open question in \cite[p.~16]{MMPS}.}
\end{figure}

Recall that $x\perpB h$ means
\[
  \|x+\lambda h\|\ge\|x\|\qquad(\lambda\in\R).
\]
An operator $T$ preserves Birkhoff--James orthogonality at $x$ when
$x\perpB h$ implies $Tx\perpB Th$ for every $h$.

\section{Counterexample}

Equip $X=\R^3$ with the norm
\begin{equation}\label{eq:norm}
 \|(a,b,c)\|=max\left\{|a|,|b|,|c|,
 \frac1{7}|-4a+2b+6c|\right\}.
\end{equation}
This is a polyhedral Banach space: its unit ball is the cube $[-1,1]^3$
intersected with the symmetric pair of slabs
$|-4a+2b+6c|\le7$.

Consider
\begin{equation}\label{eq:points}
\begin{aligned}
 x_1&=\left(1,-1,-\frac16\right),&
 x_2&=\left(\frac34,1,-1\right),&
 x_3&=\left(-1,-1,\frac56\right),\\
 y_1&=\left(1,1,-\frac56\right),&
 y_2&=(1,-1,1),&
 y_3&=\left(-1,1,\frac16\right),
\end{aligned}
\end{equation}
and the invertible operator
\begin{equation}\label{eq:T}
 T=\frac17
 \begin{pmatrix}
 4&-2&-6\\
 0&-7&0\\
 -16&-6&-25
 \end{pmatrix}.
\end{equation}

\begin{theorem}[Negative answer to the source question]\label{thm:main}
The points $x_1,x_2,x_3$ are linearly independent extreme points of
$B_X$. The operator $T$ preserves Birkhoff--James orthogonality at each of
them, but $T$ is not a scalar multiple of an isometry. Consequently the open
question in \cite[p.~16]{MMPS} has a negative answer, already in dimension
three.
\end{theorem}

The rest of the packet verifies every assertion in the theorem.

\section{An active-normal criterion}

Write
\[
 e_1=(1,0,0),\quad e_2=(0,1,0),\quad e_3=(0,0,1),\quad
 f=\frac17(-4,2,6),
\]
viewing these row vectors as functionals. The norm \eqref{eq:norm} is
$\max\{|e_1(\cdot)|,|e_2(\cdot)|,|e_3(\cdot)|,|f(\cdot)|\}$.

\begin{lemma}[Strict-sign cone test]\label{lem:cone}
Let a simple vertex $u$ of the unit ball of a three-dimensional polyhedral
norm have active norm-one supporting functionals $g_1,g_2,g_3$, and put
\[
 G_u=\begin{pmatrix}g_1\\g_2\\g_3\end{pmatrix}.
\]
Then $u\not\perpB h$ if and only if the three entries of $G_uh$ have the
same strict sign.

Suppose $S$ is invertible, $Su=v$ is another simple unit vertex, and
\begin{equation}\label{eq:criterion}
 A=G_uS^{-1}G_v^{-1}
\end{equation}
is entrywise nonnegative with no zero row. Then $S$ preserves
Birkhoff--James orthogonality at $u$.
\end{lemma}

\begin{proof}
The supporting-functional characterization of Birkhoff--James
orthogonality says
\[
 u\perpB h\quad\Longleftrightarrow\quad
 \varphi(h)=0\text{ for some }\varphi\in J(u).
\]
At a simple vertex,
$J(u)=\operatorname{conv}\{g_1,g_2,g_3\}$. Hence such a convex combination
vanishes on $h$ precisely when
\[
 \min_k g_k(h)\le0\le\max_k g_k(h).
\]
Its negation is the strict-sign assertion.

For the second claim, argue contrapositively. If
$v\not\perpB Sh$, then $z=G_vSh$ is either strictly positive or strictly
negative entrywise. Since
\[
 G_uh=G_uS^{-1}G_v^{-1}z=Az,
\]
entrywise nonnegativity and the absence of zero rows imply that $Az$ has the
same strict sign. Thus $u\not\perpB h$.
\end{proof}

\section{Exact verification}

The signed active-functional matrices at the six points in
\eqref{eq:points} are
\begin{equation}\label{eq:Gmatrices}
\begin{array}{lll}
 G_{x_1}=\begin{pmatrix}e_1\\-e_2\\-f\end{pmatrix},&
 G_{x_2}=\begin{pmatrix}e_2\\-e_3\\-f\end{pmatrix},&
 G_{x_3}=\begin{pmatrix}-e_1\\-e_2\\f\end{pmatrix},\\[16pt]
 G_{y_1}=\begin{pmatrix}e_1\\e_2\\-f\end{pmatrix},&
 G_{y_2}=\begin{pmatrix}e_1\\-e_2\\e_3\end{pmatrix},&
 G_{y_3}=\begin{pmatrix}-e_1\\e_2\\f\end{pmatrix}.
\end{array}
\end{equation}
For example, $e_1(x_1)=(-e_2)(x_1)=(-f)(x_1)=1$, while
$|e_3(x_1)|=1/6<1$. The other five checks are identical. Each displayed
triple of active functionals is linearly independent, so all six points are
simple extreme points of $B_X$.

Let $X_0$ and $Y_0$ be the matrices whose columns are the $x_i$ and $y_i$,
respectively. Direct arithmetic gives
\begin{equation}\label{eq:basicchecks}
 \det X_0=-\frac7{12},\qquad TX_0=Y_0,\qquad \det T=4.
\end{equation}
Thus the $x_i$ are linearly independent, $Tx_i=y_i$, and $T$ is invertible.
For reproducibility,
\begin{equation}\label{eq:Tinv}
 T^{-1}=
 \begin{pmatrix}
 \frac{25}{28}&-\frac1{14}&-\frac3{14}\\
 0&-1&0\\
 -\frac47&\frac27&-\frac17
 \end{pmatrix}.
\end{equation}

Substitution of \eqref{eq:Gmatrices} and \eqref{eq:Tinv} yields
\begin{equation}\label{eq:A13}
 G_{x_1}T^{-1}G_{y_1}^{-1}
 =G_{x_3}T^{-1}G_{y_3}^{-1}
 =\begin{pmatrix}
 \frac34&0&\frac14\\
 0&1&0\\
 1&0&0
 \end{pmatrix},
\end{equation}
and
\begin{equation}\label{eq:A2}
 G_{x_2}T^{-1}G_{y_2}^{-1}
 =\begin{pmatrix}
 0&1&0\\
 \frac47&\frac27&\frac17\\
 1&0&0
 \end{pmatrix}.
\end{equation}
Both matrices are entrywise nonnegative and every row is nonzero. Lemma
\ref{lem:cone}, applied three times, proves that $T$ preserves
Birkhoff--James orthogonality at $x_1,x_2,x_3$.

It remains only to rule out a scalar isometry. The point $p=(1,1,1)$ is
another unit vertex because $|f(p)|=4/7$. But
\begin{equation}\label{eq:distortion}
 Tp=\left(-\frac47,-1,-\frac{47}{7}\right),\qquad
 f(Tp)=-\frac{40}{7},\qquad
 \|Tp\|=\frac{47}{7}.
\end{equation}
On the other hand, $\|x_i\|=\|Tx_i\|=\|y_i\|=1$. If $T=cU$ for a linear
isometry $U$, any one of the $x_i$ would force $|c|=1$, contradicting
\eqref{eq:distortion}. This completes the proof of Theorem \ref{thm:main}.

\section{Independent checks and scope}

\begin{enumerate}[leftmargin=*]
\item \textbf{Full local preservation.} The proof uses the complete set
  $x_i^{\perp_B}$ through Lemma \ref{lem:cone}; it is not merely preservation
  with respect to one supporting hyperplane. This distinction is why the
  obvious diagonal examples in $\ell_1^3$ do not answer the source question.

\item \textbf{No numerical tolerance.} Every coordinate, norm, determinant,
  and cone matrix in the final counterexample is rational. The exploratory
  search was numerical, but the proof is entirely exact.

\item \textbf{Minimal claimed scope.} This disproves the universal statement
  in dimension three. It does not address whether special classes such as
  regular polygons, $\ell_1^n$, or $\ell_\infty^n$ admit positive theorems;
  the literature cited by the source contains such rigidity results.

\item \textbf{Reproducibility.} The accompanying attempt folder contains the
  polygon, cube, and general polyhedron cone-search scripts. Their only role
  was discovery and stress testing; no computation is needed to validate
  equations \eqref{eq:A13}--\eqref{eq:distortion}.
\end{enumerate}

\section{Novelty search and review recommendation}

A bounded search on August 11, 2026 used the exact question, arXiv identifier,
and combinations of ``linearly independent extreme points'', ``polyhedral
Banach space'', and ``Birkhoff--James orthogonality preservation''. It found
the source, the authors' earlier refinement paper, and positive results for
all extreme points and several classical polyhedral spaces, but no later
resolution of this exact question. The counterexample is therefore labeled a
candidate new full negative answer.

Human review should verify the strict-sign formulation of Lemma
\ref{lem:cone}, the six active-functional lists in \eqref{eq:Gmatrices}, and
the novelty search. The remaining verification is elementary matrix
arithmetic.

\begin{thebibliography}{9}
\bibitem{MMPS}
J. Manna, K. Mandal, K. Paul, and D. Sain,
\emph{An improvement of the Blanco--Koldobsky--Turn\v{s}ek characterization
of isometries}, arXiv:2511.20345 (2025).
\end{thebibliography}

\end{document}
